% Partie : théorie de la démonstration
% Chapitre : déduction naturelle
% Section : traduction de G2i vers N2i

\documentclass[../../../include/open-logic-section]{subfiles}

\begin{document}

\olfileid[fr]{pt}{nat}{tgi}

\olsection{Traduction de \Log{G2i} vers \Log{N2i}}

Nous prenons ici pour \Log{G2i} la version des règles
de \Log{G2c} dans laquelle tous les séquents ont au plus
une formule à droite, avec une adaptation de la règle
\LeftR{\lor} : ses deux prémisses et sa conclusion ont
le même conséquent, et les contextes à gauche s'additionnent.
Les antécédents de ses séquents
sont des multiensembles~$\Gamma$ de !!{formula}s, tandis
que ceux de \Log{N2i} sont des ensembles~$\Gamma'$ de
formules étiquetées, où chaque étiquette apparaît une
seule fois. Un ensemble~$\Gamma'$ de !!{formula}s
étiquetées \emph{correspond} au multiensemble~$\Gamma$
si, pour toute !!{formula}~$!A$, le nombre de ses
!!{element}s de la forme $x : !A$ est inférieur ou égal
à la multiplicité de~$!A$ dans~$\Gamma$, c'est-à-dire
au nombre de ses copies.

\begin{prop}\ollabel{prop:G2i-to-N2i}
Si $\Log{G2i} + \Cut \Proves \Gamma \Sequent \Delta$,
il existe un contexte $\Gamma'$ correspondant à~$\Gamma$
tel que $\Log{N2i} \Proves \Gamma' \Sequent !C$,
où $!C \ident \lfalse$ lorsque $\Delta$ est vide,
et $\Delta = \{!C\}$ sinon.
\end{prop}

\begin{proof}
Soit $\pi$ une !!{proof} de $\Gamma \Sequent \Delta$
dans \Log{G2i} augmenté de la coupure. Construisons
un~$\Gamma'$ correspondant à~$\Gamma$ et une
!!{proof}~$\delta$ de~$\Gamma' \Sequent !C$
dans~\Log{N2i}, par récurrence sur $n = \pheight{\pi}$.

Si $n = 0$, $\pi$ est un axiome. Pour l'axiome
$\lfalse \Sequent \quad$, $\delta$ est l'axiome
$x:\lfalse \Sequent \lfalse$ ; on prend
$\Gamma' = \{x:\lfalse\}$ et $!C \ident \lfalse$.
Pour un axiome $!A \Sequent !A$, on prend
$x: !A \Sequent !A$.

Si $n>0$, $\pi$ se termine par une règle~$R$.
L'hypothèse de récurrence donne dans~\Log{N2i}
des !!{proof}s correspondant aux prémisses.
Par renommage, on rend disjointes les étiquettes des
différentes sous-preuves. On peut aussi réserver à
chaque classe déchargée une étiquette~$x$ qui ne figure
que dans la branche au-dessus de sa décharge.
Une étiquette encore ouverte à la conclusion n'est
donc déchargée par aucune inférence de cette sous-preuve.
Les constantes propres sont renommées, au besoin,
avant toute greffe, pour préserver leur fraîcheur.
Combinons ces !!{proof}s pour obtenir un
séquent~$\Gamma' \Sequent !C$ convenable.

\begin{enumerate}
  \item Si $R$ est \RightR{\Weakening}, $\pi$ se termine par
\[
    \AxiomC{}
    \RightLabel{$\pi_1$}
    \Deduce$\Gamma \fCenter$
    \RightLabel{\RightR{\Weakening}}
    \UnaryInf$\Gamma \fCenter !A$
    \DisplayProof
    \]
  L'hypothèse de récurrence appliquée à $\pi_1$ donne une
  !!{proof} de~$\Gamma' \Sequent \lfalse$.
  L'application de $\FalseInt$ donne alors une
  !!{proof} de $\Gamma' \Sequent !A$.
  \item Si $R$ est \LeftR{\Weakening}, $\pi$ se termine par
\[
    \AxiomC{}
    \RightLabel{$\pi_1$}
    \Deduce$\Gamma \fCenter \Delta$
    \RightLabel{\LeftR{\Weakening}}
    \UnaryInf$!A, \Gamma \fCenter \Delta$
    \DisplayProof
    \]
  L'hypothèse de récurrence appliquée à $\pi_1$ donne
  une !!{proof} de~$\Gamma' \Sequent !C$, que l'on
  prend pour $\delta$. Si $\Gamma'$ correspond à
  $\Gamma$, il correspond aussi à $!A, \Gamma$ :
  le nombre de $x : !A$ dans $\Gamma'$ ne dépasse pas
  la multiplicité de~$!A$ dans $\Gamma$, donc ne
  dépasse pas non plus sa multiplicité dans
  $\Gamma \cup \{!A\}$. L'union est ici additive
  pour les multiensembles.
  \item Si $R$ est \LeftR{\Contraction}, $\pi$ se termine par
\[
    \AxiomC{}
    \RightLabel{$\pi_1$}
    \Deduce$!A, !A, \Gamma \fCenter \Delta$
    \RightLabel{\LeftR{\Contraction}}
    \UnaryInf$!A, \Gamma \fCenter \Delta$
    \DisplayProof
    \]
  L'hypothèse de récurrence appliquée à $\pi_1$ donne
  une !!{proof}~$\delta_1$ de
  $\Pi, \Gamma' \Sequent !C$, avec
  $\Pi \subseteq \{x : !A, y:!A\}$, et un contexte
  résiduel correspondant à celui qui reste après retrait
  des deux occurrences actives.
  Si $\Pi = \{x : !A, y:!A\}$, on remplace dans
  cette preuve les formules étiquetées $y:!A$ par
  $x:!A$. Ces deux classes sont ouvertes à la conclusion.
  Après le renommage préalable des classes déchargées,
  aucune inférence ne décharge une formule de la
  forme~$x:!B$ ni une formule de la classe à fusionner.
  Le renommage conserve donc toutes les règles et
  leurs contextes ; le résultat est une !!{proof}
  de~\Log{N2i}. Le contexte final, où les deux classes
  ont fusionné, correspond à l'antécédent après contraction.
  S'il y avait au plus une classe supplémentaire,
  aucune fusion n'est nécessaire.
  \item Si $R$ est $\RightR{\lif}$, $\pi$ se termine par
\[
    \AxiomC{}
    \RightLabel{$\pi_1$}
    \Deduce$!A, \Gamma \fCenter !B$
    \RightLabel{\RightR{\lif}}
    \UnaryInf$\Gamma \fCenter !A \lif !B$
    \DisplayProof
    \]
  L'hypothèse de récurrence donne une
  !!{proof}~$\delta_1$ de $x:!A, \Gamma' \Sequent !B$
  ou de $\Gamma' \Sequent !B$, où $\Gamma'$
  correspond à~$\Gamma$. On obtient $\delta$ en
  ajoutant une application de \Intro{\lif} à
  $\delta_1$, avec décharge ou décharge vacante,
  respectivement.
  \item Si $R$ est \LeftR{\lif}, $\pi$ se termine par
\[
    \AxiomC{}
    \RightLabel{$\pi_1$}
    \Deduce$\Gamma_1 \fCenter !A$
    \AxiomC{}
    \RightLabel{$\pi_2$}
    \Deduce$!B, \Gamma_2 \fCenter \Delta$
    \RightLabel{\LeftR{\lif}}
    \BinaryInf$!A \lif !B, \Gamma_1, \Gamma_2 \fCenter \Delta$
    \DisplayProof
    \]
  L'hypothèse de récurrence donne des !!{proof}s
  $\delta_1$ de $\Gamma_1' \Sequent !A$ et
  $\delta_2$ de $x:!B, \Gamma_2' \Sequent !C$
  ou de $\Gamma_2' \Sequent !C$. Dans le dernier
  cas, $\delta_2$ convient pour~$\delta$.
  Sinon, on renomme les étiquettes de~$\delta_1$
  pour qu'elles soient disjointes de celles de~$\delta_2$.
  Alors $\Gamma_1' \cup \Gamma_2'$ correspond à
  $\Gamma_1 \cup \Gamma_2$ ; l'union de droite
  additionne les multiplicités. La classe $x:!B$,
  ouverte à la conclusion de~$\delta_2$, n'est
  déchargée par aucune inférence de cette preuve.
  On remplace chaque axiome $x:!B \Sequent !B$
  par la !!{proof}~$\delta_3$
\[
    \Axiom$x: !A \lif !B \fCenter !A \lif !B$
    \AxiomC{}
    \RightLabel{$\delta_1$}
    \Deduce$\Gamma_1' \fCenter !A$
    \RightLabel{\Elim{\lif}}
    \BinaryInf$x: !A \lif !B, \Gamma_1' \fCenter !B$
    \DisplayProof
    \]
  et, dans les contextes qui dépendent de cet axiome,
  chaque occurrence de $x:!B$ par la classe
  $x:!A \lif !B$ et le contexte de sa prémisse mineure.
  La greffe donne ainsi la !!{proof}
  $\Subst{\delta_2}{\delta_3}{x:!B}$ de
  $x:!A \lif !B, \Gamma_1', \Gamma_2' \Sequent !C$.
  \item Si $R$ est \RightR{\land}, $\pi$ se termine par
\[
    \AxiomC{}
    \RightLabel{$\pi_1$}
    \Deduce$\Gamma_1 \fCenter !A$
    \AxiomC{}
    \RightLabel{$\pi_2$}
    \Deduce$\Gamma_2 \fCenter !B$
    \RightLabel{\RightR{\land}}
    \BinaryInf$\Gamma_1, \Gamma_2 \fCenter !A \land !B$
    \DisplayProof
    \]
  L'hypothèse de récurrence donne des !!{proof}s
  $\delta_1$ de $\Gamma_1' \Sequent !A$ et
  $\delta_2$ de $\Gamma_2' \Sequent !B$.
  Après renommage des étiquettes de~$\delta_2$,
  les deux preuves n'ont aucune étiquette commune.
  Alors $\Gamma_1' \cup \Gamma_2'$ correspond à
  $\Gamma_1 \cup \Gamma_2$. L'application de
  $\Intro{\land}$ à leurs séquents finaux donne
  la !!{proof}~$\delta$ recherchée.
  \item Si $R$ est $\LeftR{\land}$, $\pi$ se termine,
  par exemple, par
\[
    \AxiomC{}
    \RightLabel{$\pi_1$}
    \Deduce$!A, \Gamma \fCenter \Delta$
    \RightLabel{\LeftR{\land}}
    \UnaryInf$!A \land !B, \Gamma \fCenter \Delta$
    \DisplayProof
    \]
  L'hypothèse de récurrence donne une
  !!{proof}~$\delta_1$ de $x:!A, \Gamma' \Sequent !C$
  ou de $\Gamma' \Sequent !C$, avec $\Gamma'$
  correspondant à~$\Gamma$. Dans le dernier cas,
  on prend $\delta_1$ pour $\delta$.
  Sinon, la classe $x:!A$, ouverte à la conclusion,
  n'est déchargée par aucune inférence. On remplace
  chaque axiome $x:!A \Sequent !A$ par la
  !!{proof}~$\delta_2$
\[
    \Axiom$x: !A \land !B \fCenter !A \land !B$
    \RightLabel{\Elim{\land}}
    \UnaryInf$x: !A \land !B \fCenter !A$
    \DisplayProof
    \]
  et chaque occurrence de $x:!A$ dans les contextes
  par $x:!A \land !B$. Le résultat
  $\Subst{\delta_1}{\delta_2}{x:!A}$ est une
  !!{proof} de $x:!A \land !B, \Gamma' \Sequent !C$.
  L'autre projection se traite de la même manière.
  \item Si $R$ est \Cut, $\pi$ se termine par
\[
    \AxiomC{}
    \RightLabel{$\pi_1$}
    \Deduce$\Gamma_1 \fCenter !A$
    \AxiomC{}
    \RightLabel{$\pi_2$}
    \Deduce$!A, \Gamma_2 \fCenter \Delta$
    \RightLabel{\Cut}
    \BinaryInf$\Gamma_1, \Gamma_2 \fCenter \Delta$
    \DisplayProof
    \]
  L'hypothèse de récurrence donne des !!{proof}s
  $\delta_1$ de $\Gamma_1' \Sequent !A$ et
  $\delta_2$ de $x:!A, \Gamma_2' \Sequent !C$
  ou de $\Gamma_2' \Sequent !C$. Dans le dernier
  cas, on prend $\delta_2$ pour $\delta$.
  Sinon, on greffe $\delta_1$ à la place de chaque
  axiome $x:!A \Sequent !A$ de $\delta_2$, en
  transportant son contexte dans les nœuds dépendants.
  On obtient la !!{proof}~$\Subst{\delta_2}{\delta_1}{x:!A}$
  de $\Gamma_1', \Gamma_2' \Sequent !C$.
\end{enumerate}

Les cas non dessinés se traitent comme suit. L'introduction
à droite de la disjonction ou de l'existentielle devient
l'introduction correspondante de N2. L'introduction à
gauche de la disjonction devient son élimination : on
ajoute l'axiome de la disjonction principale, puis on
décharge les classes désignées dans les deux branches.
Les deux branches ont le même conséquent ; lorsqu'il
est vide, toutes deux se traduisent par une conclusion
absurde. L'élimination de la disjonction fournit donc
exactement la conclusion commune requise.
Pour la négation à droite, on introduit la négation
en déchargeant la classe active ; pour la négation
à gauche, on applique son élimination à la preuve
de la formule active et à l'axiome de sa négation.

Pour l'universelle à droite, on applique l'introduction
universelle ; pour l'existentielle à gauche, on ajoute
l'axiome existentiel principal et on décharge la classe
de son instance dans l'élimination existentielle.
La fraîcheur dans le séquent G2 garantit la fraîcheur
dans les contextes traduits. Pour l'universelle à gauche,
on greffe une élimination universelle issue de l'axiome
de la formule principale à la place des axiomes de
l'instance active ; le terme d'instanciation reste clos.
Dans chaque cas à gauche, une classe active inutilisée
peut simplement être ignorée.
Les contextes restants correspondent aux multiensembles
d'origine, puisque leurs unions ajoutent les nombres
de classes issues de branches aux étiquettes disjointes.
\end{proof}



















\end{document}
